\section{Radar Processing}
Every radar system operates by transmitting electromagnetic waves into an environment and 
receiving the reflected signals to detect and locate targets, (Figure \ref{fig:radar_operation}). 
The strength of the received signal is given by the equation,
\begin{equation}
P_r = \frac{P_t G_t G_r \lambda^2 \sigma}{(4\pi)^3 R^4 L}
\end{equation}
where $P_r$ is the received power, $P_t$ is the transmitted power, $G_t$ and $G_r$ are the gains
of the transmitting and receiving antennas, $\lambda$ is the wavelength, $\sigma$ is the radar
cross-section of the target, $R$ is the range to the target, and $L$ represents system losses \cite{Skolnik2001IntroRadarSystems}.

Hence, the received signal informs us about the target's radar cross-section, reflectivity and 
range, while the time delay between transmission and reception gives us just range. The 
Doppler shift of the received signal provides information about the target's radial 
velocity. If there are multiple receiving channels, the relative phase and amplitude of 
the received signals across channels can be used to estimate the bearing and elevation of 
the target.  

\begin{figure}[htb]
    \incfig[0.6]{radar_operation}
    \label{fig:radar_operation}
    \caption[Radar illustration]{A illustrative diagram of a radar detecting an object. Objects 
    are detected through the reflection of transmitted radio waves. The time delay and power of 
    the reflected signal provides information about the range and size of the object}
\end{figure}

To measure the range of targets, radar systems need to transmit the signal as pulses. Pulses are 
generally preferred over continuous constant frequency signal because the time difference between
the transmitted pulse and the received echo can be easily measured to estimate range. The 
transmitted signal is often modulated in frequency or phase. Linearly frequency modulation,
(LFM) creates pulses that involve a continuous sweep of frequencies during the pulse duration
(see Figure \ref{fig:LFM}). Advantages of such a modulation include reduced sensitivity to
frequency dependent interference, and improved range resolution due to the effect of pulse compression \cite{Richards2005FundamentalsRadarSP}.
Through pulse compression, the long duration LFM pulse can be compressed in time, resulting a 
signal that illustrates the difference between the transmitted and received signals. One pulse
compression technique for LFM signals is described as "de-chirping", which involves mixing
(term-wise multiplying) the received signal with the transmitted signal to produce a "beat"
frequency that gives range information (See Figure \ref{fig:de-chirping}). Modelling each signal
as a cosine, the mixed signal's frequencies are given by the product of cosines identity,
\begin{equation}
    \cos(\omega_1 t)\cos(\omega_2 t) = \frac{1}{2} [\cos((\omega_1-\omega_2)t) + \cos((\omega_1+\omega_2)t)]
\end{equation}
where $\omega_1$ is the transmitted signal's frequency and $\omega_2$ is the received signal's frequency.
Since both signals' frequencies increase linearly in time, the difference (beat) frequency
$\omega_1-\omega_2$ is constant and proportional to the time delay and thus the range of the
target. By applying a low-pass filter
to the mixed signal, we obtain the lower frequency component, which is the beat frequency.  

\begin{figure}[htb]
    \incfig[1]{LFM}
    \label{fig:LFM}
\end{figure}
\begin{figure}[htb]
    \incfig[1]{de-chirping}
    \label{fig:de-chirping}
    \caption[Transmitted LFM chirp, echo, and beat signal for close and far targets.]{The transmitted LFM chirp (top), delayed echo (middle), and resulting beat
    signal after mixing and low-pass filtering (bottom), shown for a close target (left, 0.5 µs
    delay) and far target (right, 2 µs delay). The far target has a higher beat frequency that the
    close target.}
\end{figure}

Since the frequency of the transmitted and received signals are linearly increasing in time, knowing
the frequency difference between the transmitted and received signals \(\Delta \omega = \omega_1-\omega_2\) allows us 
to estimate the time delay between the signals and thus the range of the target by the formula,
\begin{equation}
    R = \frac{c \cdot \Delta t}{2} = \frac
{c \cdot \Delta \omega}{2 \cdot S}
\end{equation}
where \(c\) is the speed of light, \(\Delta t\) is the time delay between the transmitted and 
received signals, \(\Delta \omega\) is the frequency difference, and \(S\) is the slope of the frequency 
modulation (\(\theta\)/\(s^2\)). 

% For radar systems where the receiver is close to the transmitter, the transmitter signal can leak
% into the receiver, creating a strong interference signal that can cover weak pulse echoes from 
% targets. To mitigate this, radar systems rapidly switch between transmission and reception modes 
% so that the transmitter isn't active when receiving. 

The compressed pulses are often pass through a low pass filter and sampled at a lower 
rate than the original signal by an Analog to Digital Converter (ADC), attenuating the frequency 
sum \(\cos(\omega_1+\omega_2)\) while preserving the useful difference frequency 
\(\cos(\omega_1-\omega_2)\).  

Exiting the analogue domain and entering the digital domain, the frequency information in the 
radar signals can be extracted by applying a Discrete Fourier Transform (DFT) to the sampled data.
At this stage the range of clear targets can be estimated by identifying the peaks in the DFT output.
At this stage many radar systems also apply a long time Fourier transform on a larger time scale 
(slow time) to extract the Doppler frequency shift, which provides information on the target's 
radial frequency. 

To improve the detectability of weak targets, radar systems often apply coherent integration 
across multiple pulses. This involves summing the complex samples across pulses for each range bin 
after the DFT. By summing the complex range bins across multiple pulses, the signal from a target
will add constructively while noise will add incoherently improving signal to noise ratio as
illustrated in Figure \ref{fig:coherent integration} \cite{Richards2010}. Phase is preserved during integration, which is 
necessary for later angle estimation stages.

\begin{figure}[htbp]
    \centering
    \incfig[1]{phasor_bins_multi_chirp}
    \caption[Coherent integration across five chirps showing constructive summation of target returns.]{Coherent integration across five chirps for each range bin. The amplitude of each arrow
    corresponds to the presence of a target at a corresponding range from the radar, and the argument
    of the arrow corresponds to the phase of the signal for that range bin.
    Most bins contain noise-like phasors with varying direction, so their coherent sum remains small.
    The highlighted bin contains a weak but consistent DC offset across chirps, so its values
    add constructively, producing a larger resultant value in the coherent-sum.}
    \label{fig:coherent integration}
\end{figure}

Constant False Alarm Rate (CFAR) detection is a key mechanism for resolving a map of noisy range data 
into discrete detections given a target false alarm rate. It does this by masking out any samples with 
magnitude below a threshold determined by the power of the samples around it (see Figure 
\ref{fig:naive-cfar}). This threshold is calculated as,
\begin{equation}
    \label{eqn:threshold}
    T = \alpha \cdot \frac{1}{N} \sum_{i=1}^{N} x_i
\end{equation}

Where $T$ is the threshold, $\alpha$ is a factor determined by the desired false alarm rate, $N$ is 
the number of training samples, and $x_i$ are the power (magnitude squared) of the training samples. 
Under the assumption that the noise in each receiver channel is complex Gaussian, the summed
power across four receivers follows a chi-squared distribution with 8 degrees of freedom.
For this distribution, the false alarm probability as a function of $\alpha$ is \cite{Richards2010},

\begin{equation}
    P_{fa}(\alpha;N) = \left( \frac{N}{N+\alpha} \right)^N
    \sum_{k=0}^{3} \binom{N+k-1}{k}
    \left( \frac{\alpha}{N+\alpha} \right)^k
\end{equation}

where $k$ is one less than the number of receivers (here $k=3$, since quadruplet samples from
four receivers are summed before thresholding). 
Figure \ref{fig:cfar_detection} illustrates how CFAR generates detections by comparing a signal to a 
threshold determined by the power of the surrounding samples.

\begin{figure}
    \incfig[1]{cfar_detection}
    \caption[CFAR detection process.]{CFAR detection process. The threshold is calculated based on 
    the power of the samples surround the test cell, and the test cell is compared to this threshold 
    to identify peaks in the signal. The scenario features three target peaks which are each detected 
    by CFAR, while the rest of the signal lies below the CFAR threshold, excluding one outlier noise 
    sample.}
    \label{fig:cfar_detection}    
\end{figure}
